\subsection{Matematiksel Model ve Algoritma}

Walrus Optimizer (WO), sürekli optimizasyon problemlerinde her bir
morsayı bir aday çözüm olarak ele alır. Bir morsanın konumu,
problemin karar değişkenlerinden oluşan bir vektördür. WO'nun amacı,
bu vektörleri yinelemeli olarak güncelleyerek amaç fonksiyonunu en
küçükleyen veya en büyükleyen çözümü bulmaktır.

Bu bölümde, en küçükleme problemi ele alınmaktadır:

\begin{equation}
    \min_{\mathbf{x}} f(\mathbf{x}),
    \qquad
    \mathbf{LB} \leq \mathbf{x} \leq \mathbf{UB},
    \label{eq:wo_problem}
\end{equation}

Burada $f(\mathbf{x})$ amaç fonksiyonunu, $\mathbf{LB}$ ve
$\mathbf{UB}$ ise karar değişkenlerinin alt ve üst sınırlarını
göstermektedir. En büyükleme problemleri, amaç fonksiyonunun işareti
değiştirilerek en küçükleme problemine dönüştürülebilir.

\subsubsection{Başlangıç popülasyonu}

Popülasyon, $N$ adet morsadan oluşur. $d$ karar değişkeni bulunan bir
problemde, $i$'inci morsanın konumu
$\mathbf{X}_i=(X_{i,1},X_{i,2},\ldots,X_{i,d})$ biçimindedir.
Başlangıç konumları, arama uzayına düzgün dağılımlı rastgele sayılar
kullanılarak yerleştirilir:

\begin{equation}
    X_{i,j}^{0}
    =
    LB_j + u_{i,j}(UB_j-LB_j),
    \qquad
    u_{i,j}\sim\mathcal{U}(0,1),
    \label{eq:wo_initialization}
\end{equation}

Burada $i=1,\ldots,N$ ve $j=1,\ldots,d$'dir. Başlangıç popülasyonunda
morsaların yüzde 90'ı yetişkin, yüzde 10'u yavru olarak kabul edilir.
Yetişkin morsalar kendi içinde eşit sayıda erkek ve dişi bireyden
oluşur. Her aday çözümün uygunluğu $f(\mathbf{X}_i)$ hesaplanır ve
en iyi çözüm $\mathbf{X}_{\mathrm{best}}$ olarak kaydedilir.

\subsubsection{Tehlike ve güvenlik sinyalleri}

WO'da sürünün davranışını seçmek için tehlike ve güvenlik sinyalleri
kullanılır. $t$ mevcut yinelemeyi, $T$ ise en fazla yineleme sayısını
göstersin. Tehlike sinyali aşağıdaki gibi tanımlanır:

\begin{align}
    \alpha_t &= 1-\frac{t}{T}, \label{eq:wo_alpha}\\
    A_t &= 2\alpha_t, \label{eq:wo_A}\\
    R_t &= 2r_1-1, \qquad r_1\sim\mathcal{U}(0,1), \label{eq:wo_R}\\
    D_t &= A_tR_t. \label{eq:wo_danger}
\end{align}

Güvenlik sinyali ise başka bir düzgün dağılımlı rastgele sayı ile
hesaplanır:

\begin{equation}
    S_t=r_2,
    \qquad
    r_2\sim\mathcal{U}(0,1).
    \label{eq:wo_safety}
\end{equation}

Yineleme ilerledikçe $\alpha_t$ azalır. WO, bu sinyalleri kullanarak
arama davranışını göç, dinlenme, besin arama ve kaçış aşamalarından
birine yönlendirir.

\subsubsection{Göç: keşif aşaması}

Tehlike sinyalinin mutlak değeri en az 1 olduğunda sürü göç davranışı
gösterir. Algoritmik açıdan bu davranış, aday çözümlerin arama
uzayında yeni bölgelere yönlendirilmesini sağlar. $m$ ve $n$, sürüden
rastgele seçilen iki farklı bireyi göstermek üzere:

\begin{align}
    \mathbf{X}_i^{t+1}
    &=
    \mathbf{X}_i^t+\mathbf{M}_t,
    \label{eq:wo_migration}\\
    \mathbf{M}_t
    &=
    \left(\mathbf{X}_m^t-\mathbf{X}_n^t\right)
    \beta_t r_3^2,
    \qquad r_3\sim\mathcal{U}(0,1),
    \label{eq:wo_migration_step}\\
    \beta_t
    &=
    1-
    \frac{1}{
        1+\exp\left(
            -10\frac{t-T/2}{T}
        \right)
    }.
    \label{eq:wo_beta}
\end{align}

Burada $\mathbf{M}_t$ göç adımını, $\beta_t$ ise bu adımın büyüklüğünü
kontrol eden katsayıyı ifade eder. $\beta_t$ yinelemeler boyunca
azalarak aramanın başlangıcında daha geniş, ilerleyen aşamalarda ise
daha sınırlı hareketler yapılmasına katkıda bulunur.

\subsubsection{Üreme: sömürü aşaması}

$|D_t|<1$ olduğunda WO üreme aşamasına geçer. Bu aşamada güvenlik
sinyali, morsaların kıyıda dinlenme davranışını mı yoksa su altında
besin arama davranışını mı izleyeceğini belirler.

\paragraph{Kıyıda dinlenme davranışı.}

$S_t\geq 0.5$ olduğunda dinlenme davranışı uygulanır. Bu davranışta
morsalar erkek, dişi ve yavru bireyler olarak ele alınır ve her grubun
konumu farklı biçimde güncellenir.

Erkek bireylerin konumları Halton dizisi kullanılarak arama uzayına
yeniden dağıtılır. Bir uygulamada bu işlem aşağıdaki biçimde
gerçekleştirilebilir:

\begin{equation}
    X_{i,j}^{t+1}
    =
    LB_j+h_{i,j}^{t}(UB_j-LB_j),
    \qquad h_{i,j}^{t}\in[0,1],
    \label{eq:wo_male_halton}
\end{equation}

Burada $h_{i,j}^{t}$, Halton dizisinin $j$'inci boyutundaki değeridir.
Bu işlem erkek bireylerin farklı bölgeleri örneklemesini ve
popülasyon çeşitliliğinin korunmasını amaçlar.

Dişi bireylerin konumları hem eş olarak kabul edilen erkek bireyden
hem de sürünün o ana kadarki en iyi çözümünden etkilenir:

\begin{equation}
\begin{aligned}
    \mathbf{F}_i^{t+1}
    ={}& \mathbf{F}_i^t
    +\alpha_t\left(\mathbf{M}_i^t-\mathbf{F}_i^t\right)\\
    &+(1-\alpha_t)
    \left(\mathbf{X}_{\mathrm{best}}^t-\mathbf{F}_i^t\right).
\end{aligned}
\label{eq:wo_female}
\end{equation}

Burada $\mathbf{F}_i^t$ dişi bireyin, $\mathbf{M}_i^t$ ise onunla
eşleştirilen erkek bireyin konumudur. $\alpha_t$ azaldıkça erkek
bireyin etkisi azalır; en iyi çözümün etkisi artar.

Yavru bireyler için önce güvenli referans konumu hesaplanır:

\begin{align}
    \mathbf{O}_i^t
    &=
    \mathbf{X}_{\mathrm{best}}^t
    +\mathbf{J}_i^t\odot\mathbf{LF},
    \label{eq:wo_safety_position}\\
    \mathbf{J}_i^{t+1}
    &=
    \left(\mathbf{O}_i^t-\mathbf{J}_i^t\right)P,
    \qquad P\sim\mathcal{U}(0,1).
    \label{eq:wo_juvenile}
\end{align}

Burada $\mathbf{J}_i^t$ yavru bireyin mevcut konumu, $P$ rastgele
rahatsızlık katsayısı, $\odot$ bileşen bazında çarpım ve $\mathbf{LF}$
Lévy dağılımına dayalı hareket vektörüdür. Lévy hareketi aşağıdaki
biçimde modellenir:

\begin{align}
    LF
    &=
    0.05\frac{x}{|y|^{1/\lambda}},
    \qquad
    x\sim\mathcal{N}(0,\sigma_x^2),
    \quad
    y\sim\mathcal{N}(0,\sigma_y^2),
    \label{eq:wo_levy}\\
    \sigma_x
    &=
    \left[
        \frac{
            \Gamma(1+\lambda)\sin(\pi\lambda/2)
        }{
            \Gamma((1+\lambda)/2)\lambda 2^{(\lambda-1)/2}
        }
    \right]^{1/\lambda},
    \qquad
    \sigma_y=1,
    \qquad
    \lambda=1.5.
    \label{eq:wo_levy_sigma}
\end{align}

$\Gamma(\cdot)$, Gamma fonksiyonudur. Lévy hareketinin uzun adımlar
üretebilmesi, yavru bireylerin çevredeki çözümlerden uzaklaşarak farklı
bölgeleri örneklemesine olanak tanır.

\paragraph{Su altında besin arama davranışı.}

$S_t<0.5$ olduğunda algoritma besin arama davranışına geçer. Bu
davranış, tehlike sinyaline göre kaçış veya toplanma olarak ikiye
ayrılır.

$|D_t|\geq 0.5$ olduğunda kaçış davranışı uygulanır:

\begin{equation}
    \mathbf{X}_i^{t+1}
    =
    \mathbf{X}_i^t R_t
    -
    \left|
        \mathbf{X}_{\mathrm{best}}^t-\mathbf{X}_i^t
    \right|r_4^2,
    \qquad r_4\sim\mathcal{U}(0,1).
    \label{eq:wo_fleeing}
\end{equation}

Mutlak değer, vektörlerde bileşen bazında alınır. Bu güncelleme,
bireyin mevcut konumunu değiştirerek sürünün bulunduğu bölgeden
uzaklaşmasını sağlar.

$|D_t|<0.5$ olduğunda toplanma davranışı uygulanır:

\begin{align}
    \mathbf{X}_i^{t+1}
    &=
    \frac{\mathbf{X}_1+\mathbf{X}_2}{2},
    \label{eq:wo_gathering}\\
    \mathbf{X}_1
    &=
    \mathbf{X}_{\mathrm{best}}^t
    -a_1b_1
    \left|
        \mathbf{X}_{\mathrm{best}}^t-\mathbf{X}_i^t
    \right|,
    \label{eq:wo_gather_x1}\\
    \mathbf{X}_2
    &=
    \mathbf{X}_{\mathrm{second}}^t
    -a_2b_2
    \left|
        \mathbf{X}_{\mathrm{second}}^t-\mathbf{X}_i^t
    \right|,
    \label{eq:wo_gather_x2}\\
    a_k
    &=
    \beta_t r_5-\beta_t,
    \qquad
    b_k=\tan(\theta_k),
    \qquad
    r_5\sim\mathcal{U}(0,1),
    \quad
    \theta_k\in[0,\pi],
    \quad k\in\{1,2\}.
    \label{eq:wo_gather_coefficients}
\end{align}

Burada $\mathbf{X}_{\mathrm{second}}^t$, o anki popülasyondan seçilen
ikinci bir bireyin konumunu temsil eder. $\mathbf{X}_1$ ve
$\mathbf{X}_2$, en iyi ve ikinci birey çevresinde oluşturulan iki
aday yönelimdir; yeni konum bu iki yönelimin ortalamasıyla elde edilir.
$\tan(\theta_k)$ fonksiyonu $\theta_k=\pi/2$ noktasında tanımsız
olduğundan, sayısal uygulamada bu noktaya çok yakın değerler
örneklenmemelidir.

\subsubsection{WO algoritmasının sözde kodu}

Algoritmanın temel karar akışı Algoritma~\ref{alg:wo} içinde
özetlenmiştir. Her yinelemede önce tehlike sinyali kontrol edilir.
Tehlike yüksekse göç uygulanır; aksi durumda güvenlik sinyali dinlenme
ve besin arama davranışları arasında seçim yapar.

\begin{algorithm}[htbp]
\caption{Walrus Optimizer (WO)}
\label{alg:wo}
\begin{algorithmic}[1]
\Require Amaç fonksiyonu $f$, alt ve üst sınırlar $\mathbf{LB}$ ve
         $\mathbf{UB}$, popülasyon büyüklüğü $N$, en fazla yineleme $T$
\Ensure En iyi çözüm $\mathbf{X}_{\mathrm{best}}$ ve amaç fonksiyonu değeri

\State $\mathbf{X}\gets$ Denklem~\eqref{eq:wo_initialization}'e göre popülasyonu başlat
\State Bireyleri yetişkin, erkek, dişi ve yavru gruplarına ayır
\State $f(\mathbf{X}_i)$ değerlerini hesapla
\State $\mathbf{X}_{\mathrm{best}}\gets$ mevcut en iyi birey
\For{$t=0$ \textbf{to} $T-1$}
    \State $D_t\gets$ Denklem~\eqref{eq:wo_danger}'e göre tehlike sinyalini hesapla
    \State $S_t\gets$ Denklem~\eqref{eq:wo_safety}'e göre güvenlik sinyalini hesapla
    \If{$|D_t|\geq 1$}
        \State Her bireyin konumunu göç denklemleri
               \eqref{eq:wo_migration}--\eqref{eq:wo_beta} ile güncelle
    \Else
        \If{$S_t\geq 0.5$}
            \State Erkek bireyleri Halton dizisi ile yeniden dağıt
            \State Dişi bireyleri Denklem~\eqref{eq:wo_female} ile güncelle
            \State Yavru bireyleri Denklem~\eqref{eq:wo_juvenile} ile güncelle
        \Else
            \If{$|D_t|\geq 0.5$}
                \State Her bireyin konumunu kaçış davranışına
                       göre, Denklem~\eqref{eq:wo_fleeing} ile güncelle
            \Else
                \State Her bireyin konumunu toplanma davranışına
                       göre, Denklem~\eqref{eq:wo_gathering} ile güncelle
            \EndIf
        \EndIf
    \EndIf
    \State Sınırları aşan konumları problem kısıtlarına göre düzelt
    \State Yeni bireylerin amaç fonksiyonu değerlerini hesapla
    \State $\mathbf{X}_{\mathrm{best}}$ değerini güncelle
\EndFor
\State \Return $\mathbf{X}_{\mathrm{best}}$
\end{algorithmic}
\end{algorithm}

\paragraph{Uygulama notu.}
Makalenin açıklama metninde kaçış davranışı Denklem~\eqref{eq:wo_fleeing},
toplanma davranışı ise Denklem~\eqref{eq:wo_gathering} ile tanımlanır.
Makaledeki sözde kodda bu iki davranışın adları ve denklem
numaraları birbiriyle ters eşleştirilmiş görünmektedir. Yukarıdaki
algoritma, bölümdeki matematiksel açıklamayı izleyerek
$|D_t|\geq 0.5$ durumunda kaçışı, diğer durumda toplanmayı uygular.

Sınır düzeltme adımı, güncellenen çözümlerin
$\mathbf{LB}\leq\mathbf{X}_i\leq\mathbf{UB}$ koşulunu sağlaması için
eklenmiştir. Denklem~\eqref{eq:wo_problem}'deki sınırlara göre
uygunluk kontrolü, uygulamada her konum güncellemesinden sonra
yapılmalıdır.